\section{Results}\label{sec:results}

\subsection{Main Results}\label{main-results}

Table~\ref{tab:main_results} summarizes task success
rates across all conditions.

\begin{table}[h]
\centering
\caption{Task success rate (\%) on OrchestraBench (70 hard tasks).}
\label{tab:main_results}
\begin{tabular}{lcc}
\toprule
\textbf{System} & \textbf{All Hard Tasks (70)} & \textbf{Multi-Phase Tasks (42)} \\
\midrule
Static-flat & 18.6 & 14.3 \\
Static-hierarchical & 52.9 & 47.6 \\
Static-debate & 64.6 & 59.5 \\
Static-oracle & 90.0 & 88.1 \\
Random switching & 55.7 & 52.4 \\
\textbf{AdaptOrch-RT} & \textbf{78.6} & \textbf{83.3} \\
\bottomrule
\end{tabular}
\end{table}

AdaptOrch-RT achieves 78.6\% overall on hard tasks, a 14-point
improvement over the best static topology (debate at 64.6\%) and within
11.4 points of the oracle. The gains concentrate on multi-phase tasks:
83.3\% versus 59.5\% for static-debate, a 23.8-point lift. On
structurally homogeneous tasks (the 28 that never change phase),
AdaptOrch-RT scores 71.4\%, close to static-debate's 71.4\% on those
same tasks. The system correctly avoids switching when the task
structure holds steady.

Random switching at 55.7\% performs worse than static-debate. Switching
at arbitrary times disrupts execution without the benefit of structural
alignment. This confirms that the timing and direction of switches
matter, not just the act of switching itself.

\subsection{Switching Behavior}\label{switching-behavior}

On the 70 hard tasks, AdaptOrch-RT triggers an average of 1.3 switches
per task. The distribution is skewed: 28 tasks receive zero switches
(the homogeneous ones), 31 receive exactly one switch, and 11 receive
two switches. No task triggers more than two switches, reflecting the
\(\delta = 0.4\) threshold and the 5-action cooldown.

The most common switching pattern is hierarchical-to-debate (19 of 42
switches), followed by flat-to-hierarchical (11 switches) and
debate-to-flat (7 switches). The remaining 5 switches span other
transitions. This aligns with the intuition that tasks often start with
decomposable planning (where hierarchy is dispatched) and shift into
iterative refinement (where debate takes over).

\subsection{Switching Precision}\label{switching-precision}

Of the 42 total switches across all tasks, 36 (85.7\%) led to successful
completion of the post-switch task segment. We evaluate this by
comparing against counterfactual continuations: for each switch, we
re-ran the task segment from the switch point under the original
topology. In 36 cases, the original topology failed the segment while
the switched-to topology succeeded. In 4 cases (9.5\%), both topologies
succeeded (the switch was unnecessary but not harmful). In 2 cases
(4.8\%), the original topology would have succeeded and the switch
introduced a failure.

The two negative switches both occurred early in execution (within the
first 8 actions) where D-I-T estimates were still noisy. Both involved
premature switches from hierarchical to debate on tasks that appeared
iterative early on (due to an initial clarification loop) but were
actually decomposable.

\subsection{When Switching Helps Most}\label{when-switching-helps-most}

The benefit of switching scales with the D-I-T variance across phases.
We quantify this as the standard deviation of the ground-truth D-I-T
vector across annotated phases within each task. Tasks with high
intra-task D-I-T variance (top quartile, \(\sigma > 0.35\)) show a
31.2-point improvement from AdaptOrch-RT over static-debate. Tasks with
moderate variance (\(0.15 < \sigma \leq 0.35\)) show a 16.4-point
improvement. Tasks with low variance (\(\sigma \leq 0.15\)) show a
2.1-point improvement, statistically indistinguishable from zero.

This pattern makes sense. If a task's structure stays roughly constant
throughout execution, picking the right topology once (as the static
oracle does) is sufficient. Switching adds value precisely when the
task's structural demands change.

\subsection{Token Overhead}\label{token-overhead}

AdaptOrch-RT consumes 8.2\% more tokens than the static-oracle baseline
on average. The overhead comes from two sources: the D-I-T monitoring
prompt (appended to each agent call, roughly 50 tokens per call) and the
switching cost (800 tokens per switch, 1.3 switches per hard task on
average).

On multi-phase tasks where switching actually fires, the overhead rises
to 11.4\%. But these are exactly the tasks where AdaptOrch-RT recovers
23.8 points of success rate. The token-per-recovered-task ratio works
out to approximately 2,400 extra tokens per additional successful task,
which is modest compared to the 30,000-50,000 tokens a typical hard task
consumes in total.

Compared to random switching (which has the same per-switch overhead but
worse outcomes), AdaptOrch-RT gets roughly 4x the benefit per token
spent on switching.
